%
% Quantifying Ocular Confounds in EEG Decoding Without an Eye Tracker
% IEEE Healthcom 2026 -- CAMERA-READY (paper #1571313658)
%
% Compile with:
%   tectonic manuscript.tex
% or:
%   pdflatex manuscript ; pdflatex manuscript
%
% Required files in the same directory:
%   fig8_leakage_angle.png
% (Bibliography is inlined below as \thebibliography, so no .bib/.bst/.bbl
%  files and no BibTeX run are required.)
%

\documentclass[conference]{IEEEtran}
\IEEEoverridecommandlockouts  % enables the corresponding-author \thanks footnote

\usepackage{cite}
\usepackage{array}
\usepackage{amsmath,amssymb,amsfonts}
\usepackage{graphicx}
\usepackage{textcomp}
\usepackage{booktabs}
\usepackage{threeparttable}

\begin{document}

\title{Quantifying Ocular Confounds in EEG Decoding Without an Eye Tracker}

\author{%
\setlength{\tabcolsep}{10pt}%
\begin{tabular}{>{\centering\arraybackslash}p{2.0in} >{\centering\arraybackslash}p{2.0in} >{\centering\arraybackslash}p{2.0in}}
Cooper Lee\newline
\textit{Illinois Mathematics and Science Academy}\newline
Aurora, IL, USA\newline
clee1@imsa.edu
&
Andrew Wong\newline
\textit{University of Michigan}\newline
Ann Arbor, MI, USA\newline
aawong@umich.edu
&
Ankit Walishetti\newline
\textit{Columbia University}\newline
New York, NY, USA\newline
aw3852@columbia.edu
\\[2.5ex]
Luke Yin\newline
\textit{Princeton University}\newline
Princeton, NJ, USA\newline
ly3569@princeton.edu
&
Adam Park\newline
\textit{Illinois Mathematics and Science Academy}\newline
Aurora, IL, USA\newline
apark1@imsa.edu
&
Yuan Yang\newline
\textit{University of Illinois Urbana-Champaign}\newline
Urbana, IL, USA\newline
yuany@illinois.edu
\\
\end{tabular}
\thanks{Corresponding author: Andrew Wong (aawong@umich.edu).}
}

\maketitle

\begin{abstract}
Eye movements, blinks, and the corneo-retinal standing potential leak into scalp
electroencephalography (EEG) and co-vary with task structure, so a classifier credited with
decoding a cognitive state may instead be partly reading the eyes. In health applications this
is not a bookkeeping matter: a communication brain--computer interface (BCI) that is covertly
gaze-driven fails precisely for the patients with oculomotor impairment it is meant to serve.
Existing defenses either remove ocular signal or model it within a single eye-tracked
recording; none audits a completed decoding dataset that has no eye-tracker. We study a
candidate auditing instrument. A frozen linear EEG-to-gaze readout, trained where synchronized
gaze exists, induces a \emph{gaze-predictive EEG subspace} (the orthogonal projection onto the
span of its weights). The Ocular Predictability Index (OPI) is the cross-validated area under
the curve (AUC) of a paradigm-matched decoder applied to that projection, minus a within-subject
label-permutation null; it measures \emph{label predictability within the projection}, not a
fraction of the original effect and not a causal attribution. We validate against ground truth in
a physically grounded EEG-plus-gaze simulator, and add the baselines and diagnostics needed to
interpret it. Against a rank-matched random projection and a bipolar electrooculography (EOG)
proxy, the gaze-trained subspace is distinguished only in the fixational regime (OPI $0.396$ vs.\
$0.149$ and $0.144$ at $0.25^\circ$); at $2^\circ$ and above all projections are equivalent within
confidence intervals because the artifact dominates. It stays near-null on neural controls (OPI $0.007$, $17\%$ significant) where
the random projection does not ($0.072$, $92\%$). Reporting reconstruction per ocular output
rather than as one aggregate reverses the reliability verdict at the fixational extreme
(horizontal-gaze $R^2=0.47$ vs.\ aggregate $0.66$). We therefore present this as a
simulation-validated proof of concept with a preliminary real-EEG illustration, not a calibrated
transferable instrument: the decisive test, training the readout on real synchronized EEG--gaze
data and transferring to independent participants, remains to be done.
\end{abstract}

\begin{IEEEkeywords}
EEG decoding, ocular artifact, eye movements, construct validity, confound auditing,
brain--computer interface, subspace projection
\end{IEEEkeywords}

\section{Introduction}

A decoding result is only as trustworthy as the signal it rests on. In EEG, a recurring
threat to construct validity is that a classifier credited with reading a cognitive
state (attention, working-memory load, target detection, the meaning of a read
word) is partly reading the eyes. The eye is an electrical dipole, cornea-positive and
retina-negative, with a standing potential of roughly $0.4$--$1.0$\,mV~\cite{malmivuo1995bioelectromagnetism};
its rotations, together with blinks and the broadband saccadic spike potential at saccade
onset~\cite{keren2010saccadic}, project to frontal and periocular electrodes with amplitudes
that dwarf cortical signals. The field of an eye movement depends on each electrode's
\emph{position} relative to the eyes: its sign reverses across the midline and its magnitude
falls with distance, so contamination is montage- and artifact-specific. Because where and when
a person looks is rarely independent of the experimental condition, this ocular activity is
\emph{condition-correlated}: it can be decoded in lockstep with the cognitive label and inflate
apparent accuracy --- the machine-learning failure mode of shortcut
learning~\cite{geirhos2020shortcut}.

The clinical stakes are concrete. A communication BCI validated on healthy participants whose
decodable signal is substantially ocular will degrade or fail in amyotrophic lateral sclerosis
and other conditions impairing oculomotor control --- the population the device targets. An EEG
biomarker of attention or workload that is partly gaze-driven tracks looking behaviour rather
than the neural state it is billed as measuring, and will not transfer across populations whose
eye-movement statistics differ. Because eyes and task are correlated by design in most
paradigms, the confound survives group averaging and is invisible in accuracy alone.

The standard defenses are incomplete. Independent component analysis and regression
\emph{remove} ocular components~\cite{jung2000removing}, but removal is imperfect and leaves a
residual whose decodable content is unknown; overweighting the saccadic spike before ICA improves
this but still targets removal, not measurement~\cite{dimigen2020optimizing}. Deconvolution
frameworks rigorously model ocular and overlap effects, but in practice within a single,
carefully co-registered, eye-tracked recording~\cite{ehinger2019unfold}.
None answers the question an author or reviewer faces about a finished result recorded with no
eye-tracker: \emph{how much of this decodable signal could be ocular?}

The concern is not hypothetical. Mostert \emph{et al.}~\cite{mostert2018eye} showed that a
visual working-memory ``representation'' could be partly explained by systematic eye movements,
and constructed a bespoke, within-experiment control. Rotaru
\emph{et al.}~\cite{rotaru2024decoding} found that EEG decoders of the spatial focus of
\emph{auditory} attention overfit to eye-gaze and per-trial fingerprints, dropping to chance
once gaze was controlled. More broadly, the field increasingly rewards auditing its methods and
benchmarks for inflated
inference~\cite{kapoor2023leakage,brookshire2024leakage,chaibubneto2018permutations}.

This paper defines such an audit and, as importantly, establishes what it can and cannot
currently support. Our contributions are: (i) a frozen linear \emph{gaze-predictive subspace}
operator and an \emph{Ocular Predictability Index} (OPI) with a within-subject permutation null
and a full-EEG ceiling; (ii) a comparison against simple alternatives --- a rank-matched random
projection, a bipolar EOG proxy, and a frontal-channel audit --- which localises the instrument's
advantage to the fixational regime and shows it has none at large amplitudes; (iii) a reliability
gate reported \emph{per ocular output}, strictly more conservative than the aggregate it replaces;
and (iv) a subspace-overlap diagnostic for the residual failure mode of absorbing genuine neural
signal, with an account of why the Haufe activation pattern is the wrong object here.

\section{Methods}

\subsection{Design and rationale}

The instrument is deliberately linear. Although a deep network is one way to learn an
EEG-to-gaze readout where synchronized eye-tracking is
available~\cite{kastrati2021eegeyenet,hollenstein2018zuco}, the object that would be frozen and
shipped is a linear readout and the orthogonal projection onto the span of its weights, so we
implement and validate exactly that. Because ground truth is available only when gaze is
recorded or simulated, the primary validation uses a forward-model simulator with known sources.
The auditor uses every channel a dataset provides; we report a full 64-channel 10--10 montage and
treat the 30 scalp-channel ERP-CORE montage~\cite{kappenman2021erpcore} as an explicitly specified subset
used to characterise how channel count affects reliability.

\subsection{The frozen gaze-predictive operator}

Let $\mathbf{X}_{\mathrm{src}}\in\mathbb{R}^{n\times p}$ be source EEG (trials $\times$
channels) for which ocular ground truth $\mathbf{Y}_{\mathrm{src}}\in\mathbb{R}^{n\times q}$
is available ($q$ outputs: horizontal and vertical gaze, a saccade channel, and a blink
indicator). A ridge readout
$\mathbf{W}=\arg\min_{\mathbf W}\lVert \mathbf{Y}_{\mathrm{src}}-\mathbf{X}_{\mathrm{src}}\mathbf{W}\rVert_F^2+\lambda\lVert\mathbf{W}\rVert_F^2$
maps EEG to ocular state. The \emph{gaze-predictive subspace} is the column space of
$\mathbf{W}$, and the locked operator is the orthogonal projection onto it,
\begin{equation}
\mathbf{P}=\mathbf{U}\mathbf{U}^{\top},\qquad \mathbf{U}=\text{left singular vectors of }\mathbf{W},
\label{eq:proj}
\end{equation}
a symmetric idempotent $p\times p$ matrix. Projecting target EEG, $\mathbf{X}\mapsto\mathbf{X}\mathbf{P}$,
retains only the component lying in the readout's span. We name this subspace by what it
demonstrably is --- the set of EEG directions from which gaze is linearly predictable --- and not
``ocular,'' because an EEG-to-gaze readout may also capture genuine cortical activity associated
with eye movements, visual attention, motor preparation, or task structure. Establishing
specificity to ocular electrical contamination would require an independent manipulation we do
not perform; this identification problem is the principal limitation of the approach
(Section~\ref{sec:limits}).

For a target with a different montage, the readout is re-fit \emph{on source data restricted to
that montage} (no periocular interpolation); the operator is frozen per montage and applied
unchanged to every target. Montage matching by shared channel names is necessary but not
sufficient for transfer: reference scheme, electrode coordinates, cap placement, hardware,
filtering, and participant anatomy all modulate the ocular scalp projection and are not
controlled by name matching.

\subsection{The Ocular Predictability Index}

A paradigm-matched cognitive decoder (a ridge linear classifier) is trained and tested under
strictly subject-wise cross-validation on the projected data $\mathbf{X}\mathbf{P}$, yielding a
macro-averaged AUC. The OPI subtracts a within-subject label-permutation null (labels shuffled
within each subject, $B=200$):
\begin{equation}
\mathrm{OPI}=\mathrm{AUC}_{\mathrm{proj}}-\overline{\mathrm{AUC}}_{\mathrm{null}},
\end{equation}
reported alongside the full-EEG AUC as a ceiling and a permutation $p$-value.

The OPI must be read for exactly what it is: \emph{how much label information survives in the
gaze-predictive projection}. It is a difference of AUCs, not a proportion, and it is not a
fraction of the original decoding effect. A high OPI does not establish that ocular artifact
\emph{caused} the observed decoding performance --- genuine gaze-correlated neural activity,
task structure, or any other signal aligned with the projection would produce the same reading.
A low OPI is likewise not a clean bill of health: if a source-trained readout fails to capture
the target's ocular artifacts, absence of detection reflects failed transfer rather than absence
of contamination.

\subsection{Baselines}

To establish whether gaze training does work that a simpler construction would not, we compute
the OPI under three comparators on identical data, differing only in the subspace projected onto: a
\emph{random} subspace matched in rank; an \emph{EOG proxy} from bipolar channel arithmetic alone
($\mathrm{F8}-\mathrm{F7}$, $(\mathrm{Fp1}+\mathrm{Fp2})/2$); and a \emph{frontal-channel audit},
the coordinate subspace of the 17 frontopolar, anterior-frontal and frontal electrodes.

\subsection{Reliability gate and bias envelope}

\emph{Reliability} is the subject-wise cross-validated gaze-reconstruction of the readout
restricted to the target montage, computed on \emph{source} data (leave-one-subject-out over
source subjects), never on the target. Because the four ocular outputs are heterogeneous, we
report them separately rather than as one aggregate: $R^2$ for the three continuous targets
(horizontal gaze, vertical gaze, saccade channel) and AUC for the binary blink indicator, for
which an $R^2$ is not interpretable. The gate is applied to the output that the audited contrast
depends on --- horizontal gaze for lateralized paradigms.

A no-eye-tracker target inherits the score of its matched source montage; where none matches, the
gate returns \emph{inconclusive}. This establishes competence \emph{within the source
distribution} only: it does not demonstrate reliability in an unseen target dataset, which would
require held-out real participants with measured gaze. Thresholds are fit on one montage grid and
evaluated on a held-out grid (Section~\ref{sec:thresh}), reported to two significant figures, and
should be treated as preliminary. The \emph{bias envelope} is measured: where ground-truth gaze
exists, we compare the auditor's OPI to the OPI computed directly from true gaze across the
saccade-amplitude range and report the signed deviation as a two-sided margin.

\subsection{The subspace-overlap diagnostic}

The instrument's central risk is that the projection absorbs genuine neural activity that
happens to be gaze-correlated, inflating the index. For a target whose cognitive decoder has
discriminative direction $\mathbf{t}$, define the overlap
\begin{equation}
s=\frac{\lVert\mathbf{P}\mathbf{t}\rVert}{\lVert\mathbf{t}\rVert}=\cos\theta,
\label{eq:overlap}
\end{equation}
the cosine of the principal angle between $\mathbf{t}$ and the gaze-predictive subspace. This is
a geometric diagnostic, not a bound: converting $s$ into decodable leakage depends on the
decoder, noise covariance, and signal amplitude, and we make no closed-form claim. We report $s$
per audit and show empirically that leakage increases with $s$ over most of its range.

A reasonable objection is that a discriminative weight vector is not a valid activation pattern,
and that the Haufe transform $\mathbf{a}=\boldsymbol{\Sigma}_{\!X}\mathbf{t}$ should be used
instead~\cite{haufe2014interpretation}. We compute both and report the comparison, because in
this setting the transform is actively counter-productive: $\boldsymbol{\Sigma}_{\!X}$ is
dominated by the very ocular variance the diagnostic is meant to be separate from, so
$\mathbf{a}$ points into the gaze-predictive subspace almost regardless of $\mathbf{t}$
(Section~\ref{sec:leak}). The object that determines what the decoder actually reads --- and
therefore what can leak --- is the discriminative direction, and that is what we use.

\subsection{Simulator and real data}

Ground-truth validation uses a forward-model simulator on a full 64-channel 10--10 montage.
Per-trial channel patterns sum literature-grounded source topographies whose amplitudes are set
as simulation parameters within reported ranges: a horizontal corneo-retinal field
($16\,\mu$V/deg), a vertical/blink frontopolar field, a focal periocular saccadic spike scaling
with saccade size~\cite{keren2010saccadic}, a frontopolar blink ($80\,\mu$V), and a lateralized
posterior N2pc-like neural component~\cite{luck1994n2pc} at $1.8\,\mu$V, on spatially correlated
background noise. An artifact
is coupled to the cognitive label with tunable strength. Each trial is a single channel-space
feature vector; this omits the temporal morphology distinguishing sustained corneo-retinal
shifts, blink waveforms, saccadic spikes, and event-related potentials, and is a substantive
simplification. The readout is trained on 20 source subjects $\times$ 200 trials at $8^\circ$
saccades; each evaluation cell is 12 subjects $\times$ 120 trials over 12 fixed-seed studies,
with ridge $\lambda{=}1$ fixed, leave-one-subject-out macro AUC, and a $B{=}200$ within-subject
label-permutation null. Uncertainties are $95\%$ confidence intervals over the 12 studies, except
in the leakage sweep of Section~\ref{sec:leak}, which uses 10 subjects over 8 studies.
Generative equations and parameters are in the released code.

For real data we use the PhysioNet EEG Motor Movement/Imagery Dataset (\texttt{eegmmidb};
all 64 channels, BCI2000, 160\,Hz; 8 subjects)~\cite{schalk2009eegmmidb,schalk2004bci2000,pollard2026physionet}. Because \texttt{eegmmidb} has no
eye-tracker, no gaze-trained operator can be built for it; the subspace is instead estimated from
the eyes-open baseline run of four \emph{source} subjects and applied to four held-out
\emph{target} subjects' motor runs. This is an \emph{illustrative proxy}, not a test of the
proposed instrument, and the sample is too small for stable subject-level estimates.

\subsection{Implementation}

The auditor, simulator, and analyses are implemented in R (CPU-only). AUC is computed by the
rank statistic and cross-checked against the \texttt{pROC} package to
$<10^{-6}$.

\section{Results}

\subsection{The operator behaves as specified}

Unit tests confirm the rank-based AUC matches \texttt{pROC} to $<10^{-6}$; the projector of
\eqref{eq:proj} is symmetric, idempotent, and correct-rank; and end-to-end the OPI is large when
the cognitive signal lies in the projection and null when it lies outside. The trained readouts
had rank 4 on the simulator and rank 2 on real \texttt{eegmmidb}.

\subsection{Sensitivity, specificity, and what the baselines show}

On the full 64-channel montage the OPI separates condition-correlated ocular contamination from
neural signal, and does so for two distinct artifact types: saccade-coupled ocular gives
OPI $=0.429\pm0.005$ with projected AUC at the full-EEG ceiling, and a condition-correlated
eye-blink artifact gives $0.198\pm0.008$, both significant in $100\%$ of studies. Negative
controls are near-null: motor imagery $0.031\pm0.008$, covert attention $0.007\pm0.012$,
label-shuffle $0.007$, neutral $-0.003$. The covert-attention case is the informative one: the
contrast is decodable from full EEG (ceiling AUC $0.79$) yet not from the projection.

The baselines qualify this substantially (Table~\ref{tab:baselines}). At the $8^\circ$ saccade
amplitude used for the headline battery, \emph{every} projection --- including a random rank-4
subspace --- recovers OPI $\approx 0.43$. That result carries no information about the merit of
gaze training: at $8^\circ$ the ocular signal ($\approx 128\,\mu$V) so dominates the neural
source ($1.8\,\mu$V) and noise ($3\,\mu$V) that almost any subspace retains it. The methods
separate only as amplitude falls. At $0.25^\circ$ --- the fixational regime that survives
eye-movement rejection, and the regime the audit is actually for --- the gaze-trained subspace
gives $0.396\pm0.015$ against $0.149\pm0.033$ for the rank-matched random projection and
$0.144\pm0.013$ for the bipolar EOG proxy.

Specificity separates the comparators differently. Against the covert-attention neural control,
the random projection returns $0.072\pm0.019$ and is significant in $92\%$ of studies --- a
false ocular attribution --- where the gaze-trained subspace returns $0.007$ ($17\%$). The EOG
proxy is likewise specific ($0.008$, $17\%$) but insensitive in the fixational regime. The
frontal-channel audit matches the gaze-trained operator on both axes in this simulator
($0.401$ at $0.25^\circ$; $-0.004$ on covert attention), achieving with rank 17 what the
gaze-trained operator achieves with rank 4. We therefore do not claim superiority over
frontal-channel auditing on this evidence. Notably, the simulator contains no frontal neural
generator, so it cannot discriminate the two; the real-data proxy below is where the
frontal approach visibly fails.

\begin{table}[!t]
\centering
\caption{OPI by projection (mean over 12 simulated studies, full 64-channel montage; 95\% CIs
$\le0.033$ throughout, given per cell in the text). Positive controls at $8^\circ$; the
fixational row is the discriminating one.}\label{tab:baselines}
\begin{threeparttable}
{\small
\begin{tabular*}{\columnwidth}{@{\extracolsep{\fill}} l cccc @{}}
\toprule
\textbf{Scenario} & \textbf{gaze} & \textbf{random} & \textbf{EOG} & \textbf{frontal} \\
 & \emph{k}=4 & \emph{k}=4 & \emph{k}=2 & \emph{k}=17 \\
\midrule
Saccade-coupled, $8^\circ$    & 0.429 & 0.428 & 0.429 & 0.429 \\
Saccade-coupled, $0.25^\circ$\tnote{a} & \textbf{0.396} & 0.149 & 0.144 & 0.401 \\
Blink-coupled                 & 0.198 & 0.199 & 0.175 & 0.195 \\
Covert attention\tnote{b}     & \textbf{0.007} & 0.072 & 0.008 & $-$0.004 \\
Motor imagery\tnote{b}        & 0.031 & 0.072 & 0.009 & 0.004 \\
\bottomrule
\end{tabular*}}
\begin{tablenotes}
\small
\item[a] Fixational regime: the only positive-control row on which the projections differ.
\item[b] Neural negative controls; \% of studies significant is $17$ (gaze), $92$ (random),
$17$ (EOG), $8$ (frontal) for covert attention.
\end{tablenotes}
\end{threeparttable}
\end{table}

The index is graded rather than binary: as the gaze--label coupling $\rho$ increases from 0 to 1
in steps of $0.2$, the OPI rises monotonically and near-linearly ($-0.013$, $0.100$, $0.202$,
$0.303$, $0.404$, $0.499$), with the projected AUC tracking from chance to $1.0$.

\subsection{Per-output reliability changes the verdict}

Aggregating reconstruction across the four heterogeneous ocular outputs masks the failure that
matters (Table~\ref{tab:peroutput}). Blink detection is near-perfect everywhere (AUC $\ge0.93$)
and vertical gaze is reconstructed well even at small amplitudes, so the aggregate stays high;
but horizontal gaze --- the dimension a lateralized paradigm depends on --- collapses at the
fixational extreme, to $R^2=0.474$ on the full montage and $0.353$ on the 30-channel subset.
Under the aggregate score the full montage at $0.25^\circ$ passes the confirmatory threshold
($0.664$); under the horizontal-gaze score it does not ($0.474$). Reporting per output is thus
strictly more conservative and reverses a verdict, which is the behaviour an audit instrument
should have.

\begin{table}[!t]
\centering
\caption{Per-output reconstruction ($R^2$ for continuous targets, AUC for the binary blink
indicator). The aggregate hides the horizontal-gaze collapse.}\label{tab:peroutput}
\begin{threeparttable}
{\small
\begin{tabular*}{\columnwidth}{@{\extracolsep{\fill}} l cccc c @{}}
\toprule
\textbf{Amp.} & \textbf{gaze$_x$} & \textbf{gaze$_y$} & \textbf{sacc.} & \textbf{blink} &
\textbf{aggr.} \\
(deg) & $R^2$ & $R^2$ & $R^2$ & AUC & $R^2$ \\
\midrule
\multicolumn{6}{l}{\emph{Full montage (64 ch)}}\\
0.25 & \textbf{0.474} & 0.880 & 0.469 & 1.000 & 0.664 \\
1.0  & 0.959 & 0.872 & 0.948 & 1.000 & 0.899 \\
12.0 & 0.999 & 0.633 & 0.988 & 0.979 & 0.755 \\
\multicolumn{6}{l}{\emph{ERP-CORE subset (30 ch)}}\\
0.25 & \textbf{0.353} & 0.761 & 0.349 & 0.993 & 0.497 \\
1.0  & 0.923 & 0.754 & 0.913 & 0.991 & 0.775 \\
12.0 & 0.999 & 0.505 & 0.988 & 0.934 & 0.679 \\
\bottomrule
\end{tabular*}}
\begin{tablenotes}
\small
\item Means over 12 studies; SDs $\le0.03$ throughout.
\end{tablenotes}
\end{threeparttable}
\end{table}

\subsection{Bias envelope and held-out thresholds}\label{sec:thresh}

Comparing the auditor's OPI to the OPI from ground-truth gaze across the saccade-amplitude
range, the bias is within $\pm0.01$ from $0.5^\circ$ upward on both montages. Only at the
fixational extreme does it grow, to $-0.034$ (full) and $-0.056$ (30-channel), and in both cases
the auditor \emph{under}-states, the conservative direction. No over-statement was observed
within the conditions swept; this envelope covers simulated horizontal saccades only and does not
generalise to real blinks, microsaccades, smooth pursuit, varying references, or other
acquisition systems.

Thresholds were previously fit and evaluated on the same grid. Refitting on the full-64 grid
alone and applying the result unchanged to the held-out 30-channel grid, the confirmatory
threshold is $R^2\approx0.66$; six of seven held-out cells are tiered confirmatory and all six
have $|\mathrm{bias}|\le0.05$, while the $0.25^\circ$ cell is correctly excluded. Under this
honest calibration the confirmatory and exploratory thresholds coincide, so the evidence supports
a \emph{two}-tier gate (confirmatory / inconclusive), not the three tiers previously reported.
The value is preliminary and specific to this simulator.

\subsection{Neural leakage tracks the overlap angle}\label{sec:leak}

Sweeping a lateralized neural source from posterior to frontal with \emph{no} ocular coupling ---
so any positive OPI is pure leakage --- leakage increases with the overlap $s$ of
\eqref{eq:overlap} (Fig.~\ref{fig:leak}, Spearman $\rho=0.95$), from $+0.006\pm0.018$ at
near-orthogonality to a maximum of $+0.143\pm0.020$. The relation is not strictly monotone: at
the most frontal source position, the highest overlap tested, leakage falls back to
$+0.119\pm0.015$, so $s$ orders the regime rather than predicting leakage exactly. The overlap
computed on the decoder's discriminative direction ranks the conditions identically to the
overlap computed on the true source topography ($\rho=1.00$), so the empirical direction is an
adequate stand-in for the generator here.

The Haufe-transformed activation pattern behaves differently, and instructively: its overlap
with the gaze-predictive subspace saturates in the range $0.92$--$0.99$ across every source
position, and correlates \emph{negatively} with leakage ($\rho=-0.54$). The reason is
structural --- the data covariance is dominated by ocular variance, so
$\boldsymbol{\Sigma}_{\!X}\mathbf{t}$ is pulled into the gaze-predictive subspace whatever
$\mathbf{t}$ is. The activation pattern is the right object for interpreting \emph{what
generated} a component; it is the wrong object for asking what a decoder \emph{reads}, which is
what governs leakage.

\begin{figure}[!t]
\centering
\includegraphics[width=\linewidth]{fig8_leakage_angle.png}
\caption{Neural leakage versus the overlap $s=\cos\theta$ between the cognitive discriminative
direction and the gaze-predictive subspace. Leakage increases with $s$ ($\rho=0.95$; not strictly
monotone at the highest overlap) and is $\approx 0$ for a posterior (near-orthogonal) source.
Bars are 95\% CIs over 8 studies.}\label{fig:leak}
\end{figure}

\subsection{An illustrative real-EEG proxy audit}

On real 64-channel \texttt{eegmmidb}, with a subspace frozen from four source subjects' eyes-open
baseline and applied to four held-out target subjects, the left-versus-right imagery contrast is
weakly decodable from full EEG (ceiling AUC $0.601$) but not from the projection (projected AUC
$0.501$, OPI $-0.023$, $p=0.998$). This is a plausible negative control rather than a measured
one: without an eye-tracker we cannot verify that left- and right-hand imagery involve no
differential gaze. Rest-versus-move is the cautionary case: the channel-proxy subspace retains
frontal \emph{neural} signal it cannot separate from ocular activity (ceiling $0.672$, projected
$0.615$, OPI $+0.114$, $p=0.002$) --- a false ocular attribution. (Macro-averaging over few
subjects raises the permutation null to $\approx0.52$, so a negative OPI means projected AUC
below that null.) With eight participants these estimates are unstable and are offered as
illustration only. Their value is negative evidence: they show what a proxy subspace does wrong,
and thereby why a readout trained on real synchronized gaze --- which this experiment does not
implement --- is the necessary next step.

\section{Discussion}\label{sec:discussion}

The construct behaves as specified in simulation: sensitive to two distinct artifact types,
near-null on neural and shuffled controls, graded in the coupling, conservative in its measured
bias, and equipped with a gate that can return ``inconclusive.'' The baselines mark out what the
evidence does \emph{not} support: gaze training earns its keep only in the fixational regime, and
a 17-channel frontal audit matches the rank-4 operator throughout this simulator. What the
comparison does establish is that subspace choice matters --- a rank-matched random projection
produces significant false ocular attributions on a purely neural contrast in $92\%$ of studies.

\emph{Relation to prior work.} The construct generalizes the within-experiment control of
Mostert \emph{et al.}~\cite{mostert2018eye} and the single-domain diagnosis of Rotaru
\emph{et al.}~\cite{rotaru2024decoding} toward a target-agnostic meter, complementing ocular
removal~\cite{jung2000removing,dimigen2020optimizing} and within-recording
deconvolution~\cite{ehinger2019unfold}. Where confound-quantification in machine learning probes
a \emph{measured} nuisance~\cite{chaibubneto2018permutations}, here the nuisance is inferred from
EEG itself --- the source of both the method's reach and its central ambiguity.

\subsection{Limitations}\label{sec:limits}

The identification problem is primary. An EEG-to-gaze readout learns the EEG directions that
predict gaze; these need not be exclusively corneo-retinal, blink, and saccadic contamination,
and may include genuine cortical activity tied to eye movements, visual attention, motor
preparation, or task structure. We therefore name the object a gaze-predictive subspace and read
the OPI as label predictability within it. Demonstrating specificity to ocular electrical
contamination requires an independent manipulation --- for example, comparing minimally- and
maximally-preprocessed versions of the same recording --- which we have not performed.

The validation is simulation-based, and the simulator's linear, spatial-only generative structure
is aligned with the linear spatial operator being tested, which limits what internal success
demonstrates. Reliability is measured on source data; transfer to unseen target participants is
assumed, not shown, and montage matching by channel name does not control reference, geometry,
hardware, or anatomy. The thresholds are preliminary, support two tiers rather than three, and
were calibrated within one simulator. The bias envelope covers simulated horizontal saccades
only. The overlap score is a geometric diagnostic whose increasing relation to decodable leakage
is empirical rather than derived, and is not strictly monotone at high overlap. The real-EEG experiment uses a channel proxy, not the proposed
gaze-trained operator, on eight participants. Finally, the index addresses only the
gaze-predictive subspace and will not detect non-ocular confounds such as repeated-stimulus
fingerprints.

Accordingly we claim no calibrated, field-deployable instrument, and propose no cross-benchmark
leaderboard on this evidence: paradigms differ too much in ocular behaviour, neural sources,
montage, reference, and preprocessing for one frozen operator to be assumed valid across them.

\emph{Next steps.} The decisive experiment remains to be run: train the readout on real
synchronized EEG--gaze recordings~\cite{kastrati2021eegeyenet,hollenstein2018zuco}, measure gaze
reconstruction in held-out real participants rather than inheriting a source score, and transfer
the frozen operator to an independent dataset. Comparing minimally- against
maximally-preprocessed versions of the same recordings would test ocular specificity directly.
Because EEGEyeNet uses cued large saccades, the sub-degree regime --- where our baselines locate
the method's only advantage --- must be anchored on reading microsaccades.

\section{Conclusion}

How much of an EEG decoding result could be the eyes rather than the mind is, for most published
datasets, unanswerable because no eye-tracker was present. We have defined a frozen,
montage-matched gaze-predictive subspace and an Ocular Predictability Index that makes the
question operational, validated its behaviour against ground truth in simulation, and --- by
adding rank-matched baselines, per-output reliability, and an overlap diagnostic --- delimited
what that validation supports. Its failure modes are now measurable rather than rhetorical; the
transferable instrument is not yet established, and real synchronized-gaze validation is the next
step rather than a completed claim.

\section*{Acknowledgment}

The authors thank the dataset providers and the anonymous reviewers, whose requests for
baselines, per-output reporting, and held-out threshold calibration materially changed the
paper's conclusions. This study used only publicly available, de-identified and synthetic data,
with no new human-subjects collection. The authors declare no competing interests and no specific
external funding. Code and result tables are released.

% Bibliography inlined (IEEEtran.bst output) so this source is self-contained,
% per the IEEEtran HOWTO recommendation for submitting .tex to external parties.
\begin{thebibliography}{10}
\providecommand{\url}[1]{#1}
\csname url@samestyle\endcsname
\providecommand{\newblock}{\relax}
\providecommand{\bibinfo}[2]{#2}
\providecommand{\BIBentrySTDinterwordspacing}{\spaceskip=0pt\relax}
\providecommand{\BIBentryALTinterwordstretchfactor}{4}
\providecommand{\BIBentryALTinterwordspacing}{\spaceskip=\fontdimen2\font plus
\BIBentryALTinterwordstretchfactor\fontdimen3\font minus
  \fontdimen4\font\relax}
\providecommand{\BIBforeignlanguage}[2]{{%
\expandafter\ifx\csname l@#1\endcsname\relax
\typeout{** WARNING: IEEEtran.bst: No hyphenation pattern has been}%
\typeout{** loaded for the language `#1'. Using the pattern for}%
\typeout{** the default language instead.}%
\else
\language=\csname l@#1\endcsname
\fi
#2}}
\providecommand{\BIBdecl}{\relax}
\BIBdecl

\bibitem{malmivuo1995bioelectromagnetism}
J.~Malmivuo and R.~Plonsey, \emph{Bioelectromagnetism: Principles and
  Applications of Bioelectric and Biomagnetic Fields}.\hskip 1em plus 0.5em
  minus 0.4em\relax New York: Oxford University Press, 1995.

\bibitem{keren2010saccadic}
A.~S. Keren, S.~Yuval-Greenberg, and L.~Y. Deouell, ``Saccadic spike potentials
  in gamma-band {EEG}: Characterization, detection and suppression,''
  \emph{NeuroImage}, vol.~49, no.~3, pp. 2248--2263, 2010.

\bibitem{geirhos2020shortcut}
R.~Geirhos, J.-H. Jacobsen, C.~Michaelis, R.~Zemel, W.~Brendel, M.~Bethge, and
  F.~A. Wichmann, ``Shortcut learning in deep neural networks,'' \emph{Nature
  Machine Intelligence}, vol.~2, no.~11, pp. 665--673, 2020.

\bibitem{jung2000removing}
T.-P. Jung, S.~Makeig, C.~Humphries, T.-W. Lee, M.~J. McKeown, V.~Iragui, and
  T.~J. Sejnowski, ``Removing electroencephalographic artifacts by blind source
  separation,'' \emph{Psychophysiology}, vol.~37, no.~2, pp. 163--178, 2000.

\bibitem{dimigen2020optimizing}
O.~Dimigen, ``Optimizing the {ICA}-based removal of ocular {EEG} artifacts from
  free viewing experiments,'' \emph{NeuroImage}, vol. 207, 2020, Art. no. 116117.

\bibitem{ehinger2019unfold}
B.~V. Ehinger and O.~Dimigen, ``Unfold: an integrated toolbox for overlap
  correction, non-linear modeling, and regression-based {EEG} analysis,''
  \emph{PeerJ}, vol.~7, 2019, Art. no. e7838.

\bibitem{mostert2018eye}
P.~Mostert, A.~M. Albers, L.~Brinkman, L.~Todorova, P.~Kok, and F.~P. de~Lange,
  ``Eye movement-related confounds in neural decoding of visual working memory
  representations,'' \emph{eNeuro}, vol.~5, no.~4, 2018, Art. no.
  ENEURO.0401-17.2018.

\bibitem{rotaru2024decoding}
I.~Rotaru, S.~Geirnaert, N.~Heintz, I.~Van~de Ryck, A.~Bertrand, and
  T.~Francart, ``What are we really decoding? {U}nveiling biases in {EEG}-based
  decoding of the spatial focus of auditory attention,'' \emph{Journal of
  Neural Engineering}, vol.~21, no.~1, 2024, Art. no. 016017.

\bibitem{kapoor2023leakage}
S.~Kapoor and A.~Narayanan, ``Leakage and the reproducibility crisis in
  machine-learning-based science,'' \emph{Patterns}, vol.~4, no.~9, 2023, Art. no.
  100804.

\bibitem{brookshire2024leakage}
G.~Brookshire, J.~Kasper, N.~M. Blauch, Y.~C. Wu, R.~Glatt, D.~A. Merrill,
  S.~Gerrol, K.~J. Yoder, C.~Quirk, and C.~Lucero, ``Data leakage in deep
  learning studies of translational {EEG},'' \emph{Frontiers in Neuroscience},
  vol.~18, 2024, Art. no. 1373515.

\bibitem{chaibubneto2018permutations}
E.~Chaibub~Neto, ``Using permutations to detect, quantify and correct for
  confounding in machine learning predictions,'' \emph{arXiv preprint
  arXiv:1805.07465}, 2018.

\bibitem{kastrati2021eegeyenet}
A.~Kastrati, M.~B. P{\l}omecka, D.~Pascual, L.~Wolf, V.~Gillioz,
  R.~Wattenhofer, and N.~Langer, ``{EEGEyeNet}: A simultaneous
  electroencephalography and eye-tracking dataset and benchmark for eye
  movement prediction,'' in \emph{Proc. Neural Inf. Process. Syst. Track on
  Datasets and Benchmarks}, vol.~1, 2021.

\bibitem{hollenstein2018zuco}
N.~Hollenstein, J.~Rotsztejn, M.~Troendle, A.~Pedroni, C.~Zhang, and N.~Langer,
  ``{ZuCo}, a simultaneous {EEG} and eye-tracking resource for natural sentence
  reading,'' \emph{Scientific Data}, vol.~5, 2018, Art. no. 180291.

\bibitem{kappenman2021erpcore}
E.~S. Kappenman, J.~L. Farrens, W.~Zhang, A.~X. Stewart, and S.~J. Luck, ``{ERP
  CORE}: An open resource for human event-related potential research,''
  \emph{NeuroImage}, vol. 225, 2021, Art. no. 117465.

\bibitem{haufe2014interpretation}
S.~Haufe, F.~Meinecke, K.~G{\"o}rgen, S.~D{\"a}hne, J.-D. Haynes, B.~Blankertz,
  and F.~Bie{\ss}mann, ``On the interpretation of weight vectors of linear
  models in multivariate neuroimaging,'' \emph{NeuroImage}, vol.~87, pp.
  96--110, 2014.

\bibitem{luck1994n2pc}
S.~J. Luck and S.~A. Hillyard, ``Electrophysiological correlates of feature
  analysis during visual search,'' \emph{Psychophysiology}, vol.~31, no.~3, pp.
  291--308, 1994.

\bibitem{schalk2009eegmmidb}
G.~Schalk, ``{EEG} Motor Movement/Imagery Dataset,'' version 1.0.0, PhysioNet,
  2009, doi: 10.13026/C28G6P.

\bibitem{schalk2004bci2000}
G.~Schalk, D.~J. McFarland, T.~Hinterberger, N.~Birbaumer, and J.~R. Wolpaw,
  ``{BCI2000}: A general-purpose brain-computer interface ({BCI}) system,''
  \emph{IEEE Transactions on Biomedical Engineering}, vol.~51, no.~6, pp.
  1034--1043, 2004.

\bibitem{pollard2026physionet}
T.~Pollard, B.~E. Moody, L.~H. Lehman, B.~J. Gow, C.~Fernandes, C.~Xie,
  A.~Johnson, R.~G. Mark, and T.~Heldt, ``{PhysioNet} as a global platform for
  biomedical research,'' \emph{Nature Health}, vol.~1, no.~8, pp. 792--795,
  2026.

\end{thebibliography}

\end{document}
